\subsection{System sygnalizacji pożarowej SSP (SAP)}
Każdy budynek otrzymuje lokalną centralę SSP z własnym zasilaczem i baterią. Centrale łączy wydzielona, kompatybilna magistrala światłowodowa producenta. Działa ona na osobnej parze włókien, bez przełączników LAN. Awaria integracji z SOC nie zatrzymuje lokalnego wykrywania ani sygnalizacji.

\begin{tabularx}{\linewidth}{@{}P{35mm}rrrrY@{}}
\toprule\textbf{Budynek} & \textbf{Czujki} & \textbf{ROP} & \textbf{Sygn.} & \textbf{I/O} & \textbf{Izolatory}\\\midrule
Terminal / CSP-0 & 48 & 6 & 6 & 2 & 2\\
Hangar / CSP-1 & 20 & 2 & 2 & 2 & 2\\
Techniczny / CSP-2 & 8 & 2 & 2 & 2 & 2\\
Portiernia / CSP-3 & 4 & 2 & 2 & 2 & 2\\\midrule
\textbf{Razem} & \textbf{80} & \textbf{12} & \textbf{12} & \textbf{8} & \textbf{8}\\\bottomrule
\end{tabularx}

\textbf{Dobór detekcji.} W terminalu rozważa się punktowe czujki dymu lub wielosensorowe; w hangarze rodzaj detekcji trzeba dobrać do wysokości, wentylacji, pyłu i spalin. Liczba 20 czujek w hangarze jest rezerwą ilościową do koncepcji, a nie potwierdzonym pokryciem. Czujki zasysające lub liniowe mogą zmienić zestawienie.

\textbf{Pętle i sygnalizacja.} Przewidziano po jednej pętli adresowalnej na budynek oraz centralę z rezerwą co najmniej jednej pętli. Orientacyjne długości pętli lokalnych: 600, 300, 150 i 100 m. Izolatory ograniczają obszar uszkodzenia zwarciowego; ich rzeczywiste rozmieszczenie wynika z architektury systemu. Liczbę sygnalizatorów należy potwierdzić pomiarem słyszalności i wymaganiami dostępności.

\textbf{Sterowania.} Osiem modułów I/O to rezerwa funkcjonalna; liczba i rodzaj kanałów wymagają macierzy sterowań. Nie zadeklarowano w ten sposób kompletnego sterowania wszystkimi 12 drzwiami. Nadzór i zwolnienie wyjść wymagają zgodnych interfejsów, dodatkowych wyjść i odpowiednich zasilaczy, jeśli wyniknie to z doboru zamknięć.

\begin{note}{Zasada nadrzędna}
Alarm pożarowy i ewakuacja wynikają z uzgodnionego scenariusza i funkcji SSP. Weryfikacja obrazu przez ochronę ani wynik AI nie mogą być warunkiem uruchomienia wymaganych działań. Panel SOC otrzymuje wyłącznie informacje o zdarzeniach SSP.
\end{note}
Zakres nie rozstrzyga obowiązku DSO ani transmisji do PSP dla rzeczywistego obiektu. Należy ustalić je według jego parametrów i przepisów; zwykłe sygnalizatory nie zastępują wymaganego DSO. Sposób monitorowania alarmów pożarowych wymaga odrębnego uzgodnienia.
\src{[2] rozporządzenie ppoż. wraz ze zmianami; [8] PN-EN 54-4 -- zakres zasilaczy. Nie deklaruje się zgodności wykonawczej całego systemu}
